\documentclass[10pt,a4paper]{amsart}
\usepackage[margin=19mm]{geometry}
\usepackage{amsmath,amssymb,mathtools}
\usepackage[scr=rsfs,cal=euler]{mathalfa}
\usepackage{xfrac}
\usepackage[unicode,hidelinks,bookmarks=false]{hyperref}
\newcommand{\ie}{i.e.\ }
\newcommand{\cf}{cf.\ }
\newcommand{\resp}{respectively\ }
\newcommand{\Subsection}{\subsection}
\providecommand{\nobreakdash}{\nobreak\hbox{-}}
\setlength{\emergencystretch}{2em}
\multlinegap0pt
\usepackage{bm}
\usepackage{bbm}
\usepackage{dsfont}
\usepackage{xspace}
\usepackage{cancel}
\usepackage{mathtools}
\usepackage{amsthm}
\makeatletter
\@ifundefined{lemma}{\newtheorem{lemma}{Lemma}}{}
\@ifundefined{proposition}{\newtheorem{proposition}{Proposition}}{}
\makeatother
\makeatletter
\newcommand{\C}{{\mathbb C}}
\expandafter\ifx\csname reviewcheckblock\endcsname\relax
\newenvironment{reviewcheckblock}{\par\smallskip\noindent{\color{reviewblue}$\bullet$ \textbf{Note:}}\par\smallskip\begingroup\color{reviewblue}}{\par\endgroup\smallskip}
\fi
\makeatother
\title[Balanced Bismut torsion-parallel manifold with constant holo]{Balanced Bismut torsion-parallel manifold with constant holomorphic sectional curvature}
\author{Qingsong Wang, Fangyang Zheng}
\date{Source: arXiv:2609.26089v1 (CC BY 4.0)}
\begin{document}
\maketitle

\noindent\emph{Source and licence.} Balanced Bismut torsion-parallel manifold with constant holomorphic sectional curvature. Version arXiv:2609.26089v1 (CC BY 4.0), distributed under the Creative Commons Attribution 4.0 International licence, as recorded in the arXiv licence metadata for this exact version and shown on the abstract page https://arxiv.org/abs/2609.26089v1. Full rights notice: licenses/arxiv-2609.26089-CC-BY-4.0.txt.\par\smallskip

\noindent\textit{Editorial scope.} Counted unit: the Proposition whose source label is \texttt{prop2} in the article (source0.tex lines 699--703, source label \texttt{prop2}), delivered together with the article's own complete original proof of it (source0.tex lines 705--870, one \texttt{proof} environment of 166 lines). No mathematics is written, repaired or re-derived: statement and proof are the article's own text line for line. Required context retained uncounted: eq:Rb (lines 233--235); eq:S (lines 243--245); eq:stabS2 (lines 270--272); lemma-DABS (lines 341--345); eq:e (lines 469--474); eq:A'A'' (lines 476--478); eq:cross (lines 480--482); lemma-min (lines 488--494); lemma-SWmin (lines 524--536); eq:SWmin (lines 529--535); eq:lemma12-min-output (lines 547--550). Every definition, display and label of those lines is retained unchanged. Omitted from this package and not redistributed: 1--232, 236--242, 246--269, 273--340, 346--468, 475--475, 479--479, 483--487, 495--523, 536--546, 551--698, 704--704, 871--1614. Nothing inside a retained excerpt is omitted or altered: every retained body is delivered exactly as the article's lines are written, so no omission receipt of any kind accompanies this package. Citations: the retained excerpts carry no citation command at all, so no bibliography entry of the article is transcribed and no reference is redistributed. Reference adaptation: none was needed and none was made; every reference inside the delivered unit resolves inside it, so no unresolved reference or citation is produced. Printed slips kept literal: no printed word, symbol, hypothesis or number of the counted statement or of its proof was altered. Layout and class: the article's own class is \texttt{amsart} with packages including amsmath, amscd, amssymb, amsthm, bbm, latexsym, amsfonts, graphicx, xcolor, needspace, xy, hyperref, geometry; the wrapper is the lane's pinned \texttt{amsart} editorial wrapper at 10pt on A4, which restates the theorem-like environments the retained text needs and adds the article's own macro definitions that the retained text needs (\texttt{\textbackslash C}), with extra wrapper packages (bm, bbm, dsfont, xspace, cancel, mathtools). The delivered numbering is the unit's own. Assets: no retained span contains a figure environment, a graphics-inclusion command, a PDF-page inclusion, a table or a TikZ picture; no image file is copied into this package, the package declares no asset dependency and no image file is redistributed with it. Licence: version arXiv:2609.26089v1 of this article is distributed under the Creative Commons Attribution 4.0 International licence, as recorded in the arXiv licence metadata for this exact version and shown on the abstract page https://arxiv.org/abs/2609.26089v1. A full rights notice is delivered at licenses/arxiv-2609.26089-CC-BY-4.0.txt. Characters: every retained excerpt is pure ASCII, so the delivered engine, XeLaTeX with its default Latin Modern text font, reports no missing character.
\begin{equation} \label{eq:Rb}
R^b_{i\bar{j}k\bar{\ell}} = \frac{c}{2} \big( \delta_{ij}\delta_{k\ell} + \delta_{i\ell}\delta_{kj} \big) - \frac{1}{2}\sum_{s=1}^n T^s_{ik} \overline{ T^s_{j\ell }} - \frac{3}{4}\sum_{s=1}^n \big( T^j_{is} \overline{ T^k_{\ell s}} + T^{\ell}_{ks} \overline{ T^i_{j s}}  \big) + \frac{1}{4}\sum_{s=1}^n \big( T^{\ell}_{is} \overline{ T^k_{j s}} + T^{j}_{ks} \overline{ T^i_{\ell s}}  \big),
\end{equation}

\begin{equation} \label{eq:S}
S_{i\bar{j}} : = \sum_{k=1}^n R^b_{i\bar{j}k\bar{k}} = \frac{c}{2}(n+1) \delta_{ij} + \frac{1}{4}\big( B_{i\bar{j}} - A_{i\bar{j}}  \big), \ \ \ \ \ \ \ \forall \ 1\leq i,j\leq n.
\end{equation}

\begin{equation} \label{eq:stabS2}
(S_{i\bar{i}} + S_{k\bar{k}}-S_{j\bar{j}})\,T^j_{ik} = 0, \ \ \ \ \ \ \ \ \ \forall \ 1\leq i,j,k\leq n.
\end{equation}

\begin{lemma} \label{lemma-DABS}
Let $(M^n,g)$ be a   balanced BTP manifold with constant Chern
holomorphic sectional curvature. For any $u,v\in V$, the curvature
derivation $D_{v\bar{u}}$ commutes with $A$, $B$, and $S$.
\end{lemma}

  \begin{equation} \label{eq:e}
\left\{ \begin{aligned} A, B, S \ \mbox{are all diagonal under} \ e; \hspace{1.35cm}\\
W = \mbox{span}\{ e_1, \ldots , e_{r_B}\}; \hspace{2.65cm}\\
S_W=\mbox{diag}(\lambda_1,\ldots , \lambda_{r_B}), \ \ \lambda_1\leq \cdots \leq \lambda_{r_B}. \end{aligned}
\right.
\end{equation}

\begin{equation} \label{eq:A'A''}
 A'_{\alpha \bar{\beta}} = \sum_{i,k\in P}T^i_{\alpha k} \overline{ T^i_{\beta k}} , \ \ \ \ A''_{\alpha \bar{\beta}} = \sum_{i\in P, \gamma \in Q}T^i_{\alpha \gamma} \overline{ T^i_{\beta \gamma}} , \ \ \ \ \ \ \ \ \forall \ \alpha , \beta \in Q.
 \end{equation}

\begin{equation} \label{eq:cross}
 R^b_{\alpha\bar{k}k\bar{\beta}} = \frac{c}{2}\delta_{\alpha \beta} +\frac{1}{2} \sum_{j\in P} T^j_{\alpha k} \overline{ T^j_{\beta k} } -\frac{3}{4} \sum_{j\in P} T^k_{\alpha j} \overline{ T^k_{\beta j} }  -\frac{3}{4} \sum_{\gamma \in Q} T^k_{\alpha \gamma } \overline{ T^k_{\beta \gamma } }.
\end{equation}

\begin{lemma} \label{lemma-min}
Let $(M^n,g)$ be a balanced BTP manifold with constant Chern holomorphic sectional curvature $c>0$ and with $B$-rank $0<r_B<n$. Then we have
\begin{equation}  \label{eq:SN}
S_N = \frac{c}{2}(m+1)I_N + \frac{1}{2}A'',
\end{equation}
where $m=n-r_B$, and $A'$ and $A''$ are defined by (\ref{eq:A'A''}). In particular, $S_N $ is positive definite.
\end{lemma}

\begin{lemma} \label{lemma-SWmin}
Let $(M^n,g)$ be a balanced BTP manifold with  constant  Chern
holomorphic sectional curvature  $c>0$ and with $B$-rank
$0<r_B<n$.   Set $m=n-r_B=\dim_{\C}N$. Then  $r_B\geq 2$, the smallest
eigenvalue of $S_W$ is simple, and
\begin{equation} \label{eq:SWmin}
 \lambda_{\min}(S_W)
 = (m+1)c+\frac{\operatorname{tr}(A'')}{m},
 \qquad
 \frac{2mc}{3}\leq\operatorname{tr}(A'')
 <\frac{r_B(r_B-1)mc}{m+2r_B}.
\end{equation}
 \end{lemma}

\begin{equation}\label{eq:lemma12-min-output}
 T^1_{\alpha j}=0
 \qquad(\alpha\in Q,\ j\in P).
\end{equation}

\begin{proposition}\label{prop2}
Let $(M^n,g)$ be a balanced BTP manifold   whose  Chern holomorphic
sectional curvature   is the positive constant $c$. Suppose that the
 $B$-rank satisfies $0<r_B<n$. Then $[W,W]_c=0$ cannot occur.
\end{proposition}

\begin{proof}
  {Set $m=n-r_B$, and choose } a local unitary frame   {$e$ satisfying
\eqref{eq:e}. Write
}\[
 {P=\{1,\ldots,r_B\},\qquad Q=\{r_B+1,\ldots,n\},
}\]
{and write $\lambda_a=S_{a\bar a}$. Since $W=\operatorname{Im}T$ } and
 $[W,W]_c=0$, we have
  \[
 {T^\alpha_{pq}=0
 \quad (1\leq p,q\leq n,\ \alpha\in Q),
 \qquad
 T^s_{ij}=0
 \quad (i,j\in P,\ 1\leq s\leq n).
}\]
  {By Lemmas \ref{lemma-min} and \ref{lemma-SWmin},
}\[
 {\lambda_\alpha\geq\frac{c}{2}(m+1)>0,\qquad
 r_B\geq2,\qquad \lambda_1<\lambda_2,
}\]
{and
}\[
 {\lambda_1=(m+1)c+\frac{\operatorname{tr}(A'')}{m},\qquad
 \frac{2mc}{3}\leq\operatorname{tr}(A'')<
 \frac{r_B(r_B-1)mc}{m+2r_B}.
}\]
{By \eqref{eq:lemma12-min-output},
$T^1_{\alpha j}=0$ for $\alpha\in Q$ and $j\in P$.
}

{Every curvature derivation } preserves $W$  and commutes with $S$ by
Lemma \ref{lemma-DABS}{. It therefore preserves the one-dimensional
minimum eigenspace $\C e_1$ } of $S_W${. Consequently,
}\[
 {R^b_{p\bar q1\bar j}=0
 \qquad (1\leq p,q\leq n,\ 2\leq j\leq r_B).
}\]
{For } $2\leq i,j\leq r_B$,  {formula \eqref{eq:Rb}, the condition
$[W,W]_c=0$, and \eqref{eq:lemma12-min-output} give
}\begin{align*}
{0
 }&{=R^b_{i\bar1 1\bar j} }\\
 &{=\frac{c}{2}\delta_{ij}
   -\frac{1}{2}\sum_sT^s_{i1}\overline{T^s_{1j}}}\\
 &{\quad-\frac{3}{4}\sum_s\left(
   T^1_{is}\overline{T^1_{js}}
   +T^j_{1s}\overline{T^i_{1s}}\right)}\\
 &{\quad+\frac{1}{4}\sum_s\left(
   T^j_{is}\overline{T^1_{1s}}
   +T^1_{1s}\overline{T^i_{js}}\right)}\\
 &{=\frac{c}{2}\delta_{ij}
   -\frac{3}{4}\sum_{\alpha\in Q}
   T^j_{\alpha1}\overline{T^i_{\alpha1}}.
}\end{align*}
  {Thus the matrix
}\[
 {F=(T^j_{\alpha1})_{\alpha\in Q,\ 2\leq j\leq r_B}
 \in M_{m,r_B-1}(\C)
}\]
  {satisfies
}\[
 {F^\ast F=\frac{2c}{3}I_{r_B-1}.
}\]
{In particular, $\operatorname{rank}F=r_B-1$, and hence
}\begin{equation}{\label{eq:rm}
 m\geq r_B-1.
}\end{equation}

{We next bound } the largest eigenvalue $\lambda_{r_B}$   {of } $S_W$.   {For
$\alpha\in Q$ } and  $j\in P$,   {Equation \eqref{eq:stabS2} gives
}\[
 {(\lambda_\alpha+\lambda_{r_B}-\lambda_j)T^j_{\alpha r_B}=0.
}\]
  {Since $\lambda_{r_B}\geq\lambda_j$ and $\lambda_\alpha>0$, it follows that
}\begin{equation}{\label{eq:prop2-max-input}
 T^j_{\alpha r_B}=0
 \qquad (\alpha\in Q,\ j\in P).
}\end{equation}
  {All components with output in $Q$ vanish; the components
$T^s_{r_Bi}$ with $i\in P$ vanish because $[W,W]_c=0$; and the components
$T^s_{r_B\alpha}$ vanish by skew-symmetry and
\eqref{eq:prop2-max-input}. Therefore
}\begin{equation}{\label{eq:prop2-Arr}
 A_{r_B\overline{r_B}}=\sum_{s,t=1}^n|T^s_{r_Bt}|^2=0.
}\end{equation}

{Set
}\[
 {a=\sum_{j\in P}\sum_{\alpha\in Q}|T^{r_B}_{\alpha j}|^2,
 \qquad
 b=\sum_{\alpha,\beta\in Q}|T^{r_B}_{\alpha\beta}|^2.
}\]
{Taking $k=r_B$ and $\beta=\alpha$ in \eqref{eq:cross}, using
\eqref{eq:prop2-max-input}, and summing over $\alpha\in Q$ gives
}\begin{equation}{\label{eq:prop2-ab}
 a+b=\frac{2mc}{3}.
}\end{equation}
{The definition of $B$ uses ordered input pairs. The $P\times P$
terms vanish, the two mixed blocks both contribute $a$, and the
$Q\times Q$ block contributes $b$. Hence
}\[
 {B_{r_B\overline{r_B}}
 =\sum_{p,q=1}^n|T^{r_B}_{pq}|^2
 =2a+b
 \leq2(a+b)=\frac{4mc}{3}.
}\]
{Together with \eqref{eq:S} and \eqref{eq:prop2-Arr}, this yields
}\begin{equation}{\label{eq:max1}
 \lambda_{r_B}
 =\frac{c}{2}(n+1)+\frac{1}{4}B_{r_B\overline{r_B}}
 \leq\frac{c}{6}(3n+3+2m).
}\end{equation}

{Define
}\[
{X=\sum_{i,j\in P}\sum_{\alpha\in Q}|T^i_{\alpha j}|^2.
}\]
{Summing the diagonal instance $\beta=\alpha$ of \eqref{eq:cross}
over } $k\in P$   {and $\alpha\in Q$, and interchanging the two $P$-indices
in one of the mixed sums, gives
}\begin{equation}{\label{eq:prop2-mixed-sum}
 0=\frac{cr_Bm}{2}+\frac{1}{2}X-\frac{3}{4}X-\frac{3}{4}\operatorname{tr}(A'')
  =\frac{cr_Bm}{2}-\frac{1}{4}X-\frac{3}{4}\operatorname{tr}(A'').
}\end{equation}
If   {$X=0$, then \eqref{eq:prop2-mixed-sum} gives
$\operatorname{tr}(A'')=2cr_Bm/3$. The strict upper bound in \eqref{eq:SWmin} would then
imply
}\[
{\frac{2}{3}<
 \frac{r_B-1}{m+2r_B},
}\]
{or equivalently
}\[
{2m+r_B+3<0,
}\]
which is   {impossible. Thus $X>0$.
}

{Choose } $i,j\in P$ and   {$\alpha\in Q$ such that
$T^i_{\alpha j}\neq0$. Equation \eqref{eq:stabS2} now gives
}\[
{\lambda_i=\lambda_j+\lambda_\alpha.
}\]
{Since $\lambda_{r_B}\geq\lambda_i$ and $\lambda_j\geq\lambda_1$, the
bounds from Lemmas \ref{lemma-min} and \ref{lemma-SWmin} imply
}\begin{align}
{\lambda_{r_B}
 }&{\geq\lambda_1+\lambda_\alpha \notag}\\
 &{\geq (m+1)c+\frac{2c}{3}
       +\frac{c}{2}(m+1)
 =\frac{c}{6}(9m+13).
 \label{eq:max2}
}\end{align}
{Combining \eqref{eq:max1} and \eqref{eq:max2} gives
}\[
{9m+13\leq3n+3+2m.
}\]
{Since $n=m+r_B$, this is
}\[
{4m+10\leq3r_B.
}\]
 On the other hand,  {\eqref{eq:rm} gives $3r_B\leq3m+3$. Thus the
assumption $[W,W]_c=0$ forces $m+7\leq0$, whereas
$m=n-r_B=\dim_{\C}N\geq1$. This contradiction proves that
$[W,W]_c=0$ cannot occur.}
\end{proof}

\end{document}
